%!TEX ROOT = thesis.tex
\chapter{SYSTEM DESIGN AND ARCHITECTURE}
\label{chap:architecture}

This chapter provides an exhaustive architectural and mathematical deconstruction of the proposed Nano-Mamba U-Net. Moving beyond high-level topological descriptions, this section rigorously defines the tensor transformations, the hierarchical feature extraction pathways, and the fundamental micro-calculus underlying the Spatio-Temporal Selective State Space bottleneck. The objective is to transparently justify how the architecture achieves superior spatio-temporal modeling of cardiac dynamics within an ultra-lightweight parameter envelope of merely 1.45 million parameters.

\section{Overall Topological Architecture}
The Nano-Mamba U-Net adheres to a symmetric Encoder-Bottleneck-Decoder topology, fundamentally inspired by the standard 3D U-Net \cite{cicek20163d}, but critically deviates in its deepest representational layer. The architecture leverages consecutive 3D convolutional blocks in the contracting path to extract localized spatial hierarchies, and utilizes transposed convolutions in the expanding path for precise anatomical localization. The core innovation—the Spatio-Temporal Mamba Bottleneck—is strategically positioned at the apex of the contracting path, where the spatial dimensions are maximally compressed, and the semantic channel capacity is at its peak.

To visually articulate the complex forward propagation and information fusion mechanisms, Figure \ref{fig:full_architecture} presents the complete topological schematic of the proposed network.

\begin{figure}[hbt!]
\centering
\resizebox{0.95\textwidth}{!}{
\begin{tikzpicture}[
    node distance=1.2cm and 1.5cm,
    layer/.style={rectangle, draw=black!80, thick, align=center, rounded corners=2pt, minimum height=0.8cm},
    enc/.style={layer, fill=blue!10},
    dec/.style={layer, fill=green!10},
    mamba/.style={layer, fill=red!10, draw=red, very thick},
    arrow/.style={->, >=stealth, thick},
    skip/.style={->, >=stealth, thick, dashed, draw=blue}
]

% Encoder Path (Left Side)
\node[enc] (in) {\textbf{Input 3D MRI}\\ $1 \times 256 \times 256 \times 16$};
\node[enc, below=of in] (e1) {\textbf{Encoder 1}\\ Conv3D + ReLU\\ $16 \times 256 \times 256 \times 16$};
\node[enc, below=of e1] (e2) {\textbf{Encoder 2}\\ MaxPool $\downarrow$ + Conv3D\\ $32 \times 128 \times 128 \times 8$};
\node[enc, below=of e2] (e3) {\textbf{Encoder 3}\\ MaxPool $\downarrow$ + Conv3D\\ $64 \times 64 \times 64 \times 4$};

% Bottleneck (Bottom Center)
\node[mamba, below right=1cm and 1cm of e3] (bot) {\textbf{Spatio-Temporal Mamba Bottleneck (128 Channels)}\\ Tensor Serialization $\rightarrow$ Selective Scan $\rightarrow$ Volumetric Reshaping\\ $\mathbb{R}^{128 \times 32 \times 32 \times 2}$};

% Decoder Path (Right Side)
\node[dec, above right=1cm and 1cm of bot] (d3) {\textbf{Decoder 3}\\ UpConv $\uparrow$ + Concat + Conv3D\\ $64 \times 64 \times 64 \times 4$};
\node[dec, above=of d3] (d2) {\textbf{Decoder 2}\\ UpConv $\uparrow$ + Concat + Conv3D\\ $32 \times 128 \times 128 \times 8$};
\node[dec, above=of d2] (d1) {\textbf{Decoder 1}\\ UpConv $\uparrow$ + Concat + Conv3D\\ $16 \times 256 \times 256 \times 16$};
\node[dec, above=of d1] (out) {\textbf{Output Mask}\\ $1 \times 1 \times 1$ Conv3D\\ $4 \times 256 \times 256 \times 16$};

% Connections
\draw[arrow] (in) -- (e1);
\draw[arrow] (e1) -- (e2);
\draw[arrow] (e2) -- (e3);
\draw[arrow] (e3) |- (bot);
\draw[arrow] (bot) -| (d3);
\draw[arrow] (d3) -- (d2);
\draw[arrow] (d2) -- (d1);
\draw[arrow] (d1) -- (out);

% Skip Connections
\draw[skip] (e3) -- node[above, font=\scriptsize] {Skip Connection (Concat)} (d3);
\draw[skip] (e2) -- node[above, font=\scriptsize] {Skip Connection (Concat)} (d2);
\draw[skip] (e1) -- node[above, font=\scriptsize] {Skip Connection (Concat)} (d1);

\end{tikzpicture}
}
\caption{The complete architectural topology of the Nano-Mamba U-Net. Blue blocks represent the contracting spatial encoder, green blocks represent the expanding decoder, and the central red block denotes the linear-time Selective State Space Bottleneck. Dashed blue lines indicate spatial skip connections preserving high-frequency anatomical boundaries.}
\label{fig:full_architecture}
\end{figure}

\section{The Spatio-Temporal Mamba Bottleneck}
The empirical superiority of the proposed architecture relies entirely on the successful mathematical integration of the Selective State Space Model \cite{gu2023mamba} into the highly compressed latent space.

\subsection{Volumetric Tensor Serialization}
Unlike localized 3D convolutions that naturally process grid-like spatial structures, the Mamba architecture fundamentally requires a 1D serialized sequence to perform its hardware-aware parallel scan. Let the highly compressed spatial feature map outputted by Encoder 3 and passed through the initial bottleneck convolution be denoted as a 5-dimensional continuous tensor:
\begin{equation}
    \mathcal{X} \in \mathbb{R}^{B \times C \times H \times W \times D},
\end{equation}
where $B$ represents the mini-batch dimension, $C=128$ is the optimal semantic channel depth (as definitively proven by our subsequent ablation studies), and $H, W, D$ are the spatial coordinates ($32 \times 32 \times 2$). 

To ingest this volumetric data into the State Space ODEs, the tensor is systematically flattened across its three spatial dimensions into a continuous 1D sequence of length $L$, where $L = H \times W \times D$. This mathematical transformation is expressed as:
\begin{equation}
    \mathbf{x}_{seq} = \text{Reshape}(\mathcal{X}) \in \mathbb{R}^{B \times C \times L}.
\end{equation}
To strictly align with the causal, autoregressive sequence processing requirements of the Mamba underlying CUDA kernels, the tensor must be transposed to treat the sequence length as the primary operational dimension:
\begin{equation}
    \mathbf{x}' = \text{Transpose}(\mathbf{x}_{seq}) \in \mathbb{R}^{B \times L \times C}.
\end{equation}
This serialization essentially forces the model to traverse the 3D cardiac volume as a temporal sequence, sequentially mapping the spatial voxels across the temporal cardiac cycle.

\subsection{Continuous-Time State Space ODEs}
Once serialized, the sequence is processed by the State Space mechanism. Fundamentally, Mamba maps the 1D input sequence $x(t) \in \mathbb{R}$ to an output sequence $y(t) \in \mathbb{R}$ through an implicit, high-dimensional hidden latent state $h(t) \in \mathbb{R}^N$. This continuous-time dynamic mapping is governed by a set of linear Ordinary Differential Equations (ODEs):
\begin{equation}
    h'(t) = \mathbf{A}h(t) + \mathbf{B}x(t),
\end{equation}
\begin{equation}
    y(t) = \mathbf{C}h(t).
\end{equation}
Here, $\mathbf{A} \in \mathbb{R}^{N \times N}$ is the critically important state transition matrix that dictates the memory decay of the system (i.e., how long the network remembers the shape of the Right Ventricle from previous frames). $\mathbf{B} \in \mathbb{R}^{N \times 1}$ and $\mathbf{C} \in \mathbb{R}^{1 \times N}$ are learnable projection matrices for the input and output, respectively.

\subsection{Zero-Order Hold Discretization}
Because modern deep learning hardware operates on discrete, quantized tensors rather than continuous differential functions, the continuous ODEs (Equations 4.4 and 4.5) must be analytically discretized. Mamba achieves this through the Zero-Order Hold (ZOH) assumption, introducing a hardware-aware temporal step size parameter $\mathbf{\Delta}$. 

Applying the ZOH integral over a discrete interval $k$, the continuous matrices $\mathbf{A}$ and $\mathbf{B}$ are rigorously transformed into their discrete counterparts $\mathbf{\bar{A}}$ and $\mathbf{\bar{B}}$:
\begin{equation}
    \mathbf{\bar{A}} = \exp(\mathbf{\Delta} \mathbf{A}),
\end{equation}
\begin{equation}
    \mathbf{\bar{B}} = (\mathbf{\Delta} \mathbf{A})^{-1} (\exp(\mathbf{\Delta} \mathbf{A}) - \mathbf{I}) \cdot \mathbf{\Delta} \mathbf{B}.
\end{equation}
Following this discretization, the continuous differential system collapses into an elegant, discrete linear recurrence relation:
\begin{equation}
    h_k = \mathbf{\bar{A}}h_{k-1} + \mathbf{\bar{B}}x_k,
\end{equation}
\begin{equation}
    y_k = \mathbf{C}h_k.
\end{equation}

\subsection{The Selective Mechanism and Linear Complexity}
In standard State Space Models, matrices $\mathbf{\bar{A}}$, $\mathbf{\bar{B}}$, and $\mathbf{C}$ are invariant across time, which heavily limits their ability to process dynamic, context-dependent inputs. The profound architectural innovation of the Selective State Space Model is making the step size $\mathbf{\Delta}$, as well as matrices $\mathbf{B}$ and $\mathbf{C}$, strictly data-dependent functions of the input $x_k$. 

This "Selective Scan" mechanism endows the network with the mathematical capability to dynamically filter information: it can actively selectively remember critical structural features (e.g., the sharp boundary of the myocardium) and actively forget irrelevant artifacts (e.g., MRI bias field noise). Because Equation 4.8 is a pure recurrence relation that updates the hidden state $h_k$ without needing to compute the $L \times L$ pairwise interactions demanded by Vision Transformers, the computational complexity drops from quadratic $\mathcal{O}(L^2)$ to strictly linear $\mathcal{O}(L)$. 

Finally, the processed 1D output sequence $y_k$ is transposed and reshaped back into its original volumetric shape $\mathbb{R}^{B \times 128 \times 32 \times 32 \times 2}$ before being passed to the decoder pathway for high-resolution anatomical reconstruction. This elegant combination of spatial convolutions and linear-time selective state sequences yields the unprecedented 95.83\% accuracy within a minuscule 1.45M parameter envelope.

% =========================================================
% 下面这两节是咱们新加的微观算法剖析与超参数配置！
% =========================================================

\section{Micro-Architectural Components}
The efficacy of the Nano-Mamba U-Net is not merely a function of its global topology but is fundamentally rooted in the mathematical properties of its internal processing units. This section deconstructs the localized spatial feature extraction blocks and the specific configuration of the state space equations.

\subsection{Standard 3D Double-Convolutional Block}
Each stage of the encoder and decoder pathways utilizes a "DoubleConv" block, comprising two consecutive $3 \times 3 \times 3$ convolutional layers. Mathematically, these layers are augmented with 3D Batch Normalization (BN) \cite{ioffe2015batch} and Rectified Linear Unit (ReLU) activation functions. 

The integration of 3D BN is critical for stabilizing the internal covariate shift when processing heterogeneous MRI volumes. For a mini-batch of feature maps $\mathcal{B}$, the normalization for a specific channel is formulated as:
\begin{equation}
    \hat{x} = \frac{x - E[\mathcal{B}]}{\sqrt{Var[\mathcal{B}] + \epsilon}},
\end{equation}
followed by a learnable affine transformation $y = \gamma \hat{x} + \beta$. This ensures that the gradient flow remains within a stable numerical range, preventing the vanishing gradient problem in deep 3D networks. Following normalization, the non-linearity is introduced via the ReLU function:
\begin{equation}
    f(x) = \max(0, x).
\end{equation}
By maintaining a constant gradient of 1 for all positive activations, ReLU facilitates sparse representations and accelerates the convergence of the Adam optimizer.

\subsection{Selective Scan Algorithm and Implementation}
The core logic of the Spatio-Temporal Mamba Bottleneck is the Selective Scan algorithm. Unlike standard RNNs that suffer from sequential processing bottlenecks, Mamba utilizes a parallelizable prefix sum approach to compute the recurrent states. The algorithmic implementation of this process, adapted for 3D tensor processing in this research, is detailed in Algorithm \ref{alg:mamba}.

\begin{algorithm}[hbt!]
\caption{Selective Scan Mechanism for 3D Cardiac Tensors}
\label{alg:mamba}
\begin{algorithmic}[1]
\REQUIRE Input serialized tensor $\mathbf{x} \in \mathbb{R}^{B \times L \times C}$, Discretization parameter $\mathbf{\Delta}$
\ENSURE Transformed spatio-temporal features $\mathbf{y} \in \mathbb{R}^{B \times L \times C}$
\STATE \textbf{Step 1: Data-Dependent Projection}
\STATE $\mathbf{B}, \mathbf{C} \leftarrow \text{Linear}(\mathbf{x})$, $\mathbf{\Delta} \leftarrow \text{Softplus}(\text{Linear}(\mathbf{x}))$
\STATE \textbf{Step 2: Discretization via ZOH}
\STATE $\mathbf{\bar{A}} = \exp(\mathbf{\Delta} \mathbf{A})$
\STATE $\mathbf{\bar{B}} = (\mathbf{\Delta} \mathbf{A})^{-1} (\exp(\mathbf{\Delta} \mathbf{A}) - \mathbf{I}) \cdot \mathbf{\Delta} \mathbf{B}$
\STATE \textbf{Step 3: Recurrent Parallel Scan}
\FOR{$t = 1$ \TO $L$}
    \STATE $h_t = \mathbf{\bar{A}}_t h_{t-1} + \mathbf{\bar{B}}_t x_t$
    \STATE $y_t = \mathbf{C}_t h_t$
\ENDFOR
\RETURN $\mathbf{y}$
\end{algorithmic}
\end{algorithm}

\section{Implementation Details and Hyperparameters}
To ensure the reproducibility of the Nano-Mamba U-Net's 95.83\% DSC performance, the training environment and hyperparameter configurations were strictly standardized. All weights were initialized using the Kaiming Uniform distribution \cite{he2015delving} to maintain a stable variance of activations across layers.

\begin{table}[hbt!]
\caption{Global Hyperparameter Configuration for Nano-Mamba U-Net Training.}
\label{tab:hyperparams}
\centering
\begin{tabular}{l | p{8cm}}
\hline
\textbf{Hyperparameter} & \textbf{Selection Rationale and Value} \\
\hline
Initial Learning Rate & $1 \times 10^{-4}$ (Selected to ensure stable convergence in the early epochs). \\
Batch Size & 1 (Constrained by the volumetric 3D tensor dimensions and 8GB VRAM limit). \\
Optimizer & Adam ($\beta_1=0.9, \beta_2=0.999$, $\epsilon=1e-8$). \\
Weight Decay & $1 \times 10^{-5}$ (Implemented to prevent spatial overfitting in the 100-patient cohort). \\
Epochs & 150 (Total iterations required for loss stabilization). \\
Data Augmentation & Random 3D Rotation ($\pm 15^\circ$), Scaling ($[0.9, 1.1]$), and Gaussian Noise. \\
\hline
\end{tabular}
\end{table}

The choice of a mini-batch size of 1 is a strategic clinical constraint. In point-of-care environments, data is often processed patient-by-patient; thus, optimizing the model for singular volumetric inference ensures its immediate readiness for edge-device deployment.
